\chapter*{M1: Memory Allocator}{\textbf{Author: Chenyi Wang}
\newline
 \newline
 The goal of this milestone is to enable a privileged process to distrubute available
 physical memory to other processes.}
\addcontentsline{toc}{chapter}{M1: Memory Allocator}
\section{Architecture Overview}
\subsection{Server Placement}

We implemented the memory server inside the init domain rather than a separate
process. The init domain owns the only mm instance, \co{aos_mm}, defined as a
file-static variable in \co{usr/init/mem_alloc.c}. Every other domain requests
RAM by sending an \co{AOS_RPC_MSG_RAM_REQ} over its LMP channel to init. The
handler \co{handle_ram_req} in \co{usr/init/rpc_server.c} calls
\co{aos_ram_alloc_aligned}, which is a thin wrapper around \co{mm_alloc_aligned},
and replies with the resulting RAM cap.

\medskip\noindent
This kept the bootstrap problem trivial. The init endpoint cap is already in
every child's CSpace as \co{cap_initep}, so no name service or extra cap exchange
is needed. We can move the memory server into its own domain later without
changing the client side, since \co{aos_rpc_get_memory_channel} would just
return a different \co{aos_rpc}.

\subsection{Disjoint Allocations}

There is exactly one mm instance, and all requests are serialised through the
single init dispatcher thread, so two requests can never be served from
overlapping regions. \co{mm_alloc_aligned} reserves a region from the free list
before it returns, and the capability machinery prevents the same physical
memory from being retyped to RAM twice. Disjointness is a structural property
of the design; the handler does not need to check it.

\subsection{Per-Client Quota}

To stop one misbehaving client from grabbing all of physical memory, every
\co{client_state} carries a \co{bytes_allocated} counter. The handler increments
it after each successful allocation. A request is rejected with
\co{LIB_ERR_RAM_ALLOC} when \co{bytes_allocated} plus the requested size would
exceed the quota, currently 256 MiB. The quota is hard-coded in
\co{handle_ram_req} but easy to move to a configurable constant later.

\medskip\noindent
One important caveat is that the counter only tracks allocations, not frees. Once we
add an explicit free RPC, the counter has to be decremented.

\subsection{Alignment}

The alignment value the client passes in \co{aos_rpc_get_ram_cap} goes through
to \co{aos_ram_alloc_aligned} untouched. We do not validate or round it inside
the handler; the mm library is the authority on what alignments are achievable.

\subsection{Error Handling}

A failure reply carries a \co{NULL_CAP} cap, \co{actual_bytes} of 0, and an
\co{errval_t} in \co{words[2]}. The client-side \co{aos_rpc_get_ram_cap} reads
\co{words[2]} as the error code and returns immediately on failure, so the
caller never sees a partially filled \co{capref}.

\subsection{Switching Children to the Remote Allocator}

Once a child establishes its RPC channel to init inside
\co{barrelfish_init_onthread}, we call \co{ram_alloc_set(NULL)}.
This clears the local allocator function pointer. After that,
\co{ram_alloc_aligned} falls through \co{ram_alloc_fixed} and into
\co{ram_alloc_remote} once the early boot RAM buffer is exhausted, and
\co{ram_alloc_remote} calls \co{aos_rpc_get_ram_cap}. The effect is that
\co{frame_alloc} and \co{malloc} transparently route through the memory server.
The init domain returns early from \co{barrelfish_init_onthread} via the
\co{if (init_domain) return} guard, so init never calls \co{ram_alloc_set(NULL)}
and there is no risk of it calling its own RPC server.

\subsection{Limitations}

We have not implemented a per-client free path yet, so caps returned to a
client stay accounted against its quota for the lifetime of the system.
The 256 MiB quota is
not tuned for the actual physical memory size of the target board; on a
1~GiB iMX8X system this allows at most four clients before quota errors
kick in.

\section{Development Process}

For this milestone, we split the parts evenly between each other. First, Allan
implemented the \co{lib/mm} library, including support for partial frees and 
range allocations. For mapping frame capabilities through a linear allocator, Sahil implemented the
main paging code which enabled this feature. At this point, we did misunderstand
some of the specification, as we thought we only required one L3 page table for 
this milestone.  Shane then finished up the pre-testing code for the milestone by implementing the 
slab allocator refill code.

\medskip\noindent
In terms of debugging most of the fault was with the paging code as we were under the impression
we only needed a single L3 leaf page table. Once we fixed this using the same linear allocation
pattern, the rest of the fixes we made were quite straightforward and we were able to finish the
milestone.
